% !PRINTED 625
\weberpage{625}

Da hier $b=0$ ist, so wird nach (7) und (8) $y=2u$, und
aus (11) und (15) ergeben sich die Werthe der $y$:
\[
\begin{aligned}
y_1&=y_2=y_3=y_6=-2\sqrt{\alpha}\\
y_4&=(4-2i)\sqrt{\alpha},\quad y_5=(4+2i)\sqrt{\alpha}.
\end{aligned}
\]
Danach erh\"alt man f\"ur $\beta=0$ die folgende in $y$ identische
Gleichung:
\weq{17}{\begin{aligned}
&y^6+a_2y^4+a_4y^2+a_6-a_5\sqrt{\Delta}\,y\\
&=(y+2\sqrt{\alpha})^4(y^2-8y\sqrt{\alpha}+20\alpha)\\
&=y^6-20\alpha y^4+240\alpha^2y^2+512\sqrt{\alpha}^5y+320\alpha^3.
\end{aligned}}
Daraus findet man aber die allgemeine Resolvente f\"ur die
Gleichung (13), wenn man nach (16) f\"ur $-16\sqrt{\alpha}^5$ setzt $\sqrt{\Delta}$, also
\weq{18}{y^6-20\alpha y^4+240\alpha^2y^2-32\sqrt{\Delta}\,y+320\alpha^3=0.}

Die Discriminante $\Delta$ hat hier, wie sich aus der Formel (3)
im \S.\,74 ergiebt, den Ausdruck
\weq{19}{\Delta=5^5\beta^4+2^8\alpha^5.}

Einfacher noch wird die Gleichung f\"ur $u$, n\"amlich
\weq{20}{u^6-5\alpha u^4+15\alpha^2u^2-\sqrt{\Delta}\,u+5\alpha^3=0,}
und wenn man $u^2=v$ setzt, so ergiebt sich f\"ur $v$ die rationale
Gleichung 6$^{\text{ten}}$ Grades
\weq{21}{(v^3-5\alpha v^2+15\alpha^2v+5\alpha^3)^2=\Delta v.}

Eine andere Form dieser Gleichung erh\"alt man aus (17),
wenn man den zweiten der drei Ausdr\"ucke mit dem multiplicirt,
den man daraus erh\"alt, wenn man $\sqrt{\alpha}$ in $-\sqrt{\alpha}$ verwandelt.
Dadurch bekommt man
\[
(v-\alpha)^4(v^2+6\alpha v+25\alpha^2)=0,
\]
und in diesen muss (21) \"ubergehen, wenn $\beta=0$ gesetzt wird.
Da aber in (21) der von $\beta$ abh\"angige Theil durch (19) bestimmt
ist, so folgt die gesuchte Form der Resolvente
\weq{22}{(v-\alpha)^4(v^2+6\alpha v+25\alpha^2)=5^5\beta^4v.}

Man kann noch die Frage aufwerfen, ob die hier eingef\"uhrte
Function $v$ immer wirklich metacyklisch ist, ob sie nicht vielleicht
 bei besonderen Gleichungen noch bei anderen Permutationen
 unge\"andert bleibt.

Wenn dieser Fall eintreten sollte, so m\"usste die Gleichung
(21) oder (22) gleiche Wurzeln oder (20) gleiche oder entgegengesetzte
 Wurzeln haben.
